Las \textit{Blue Stragglers} son estrellas que aparecen más calientes y luminosas que el \textit{turn-off point} de su cúmulo estelar, aparentando haber ``rejuvenecido''. Para explicar su origen existen dos hipótesis: la \textit{hipótesis aislada} y la \textit{hipótesis de transferencia de masa binaria}. Ambos escenarios producen estrellas en idénticas posiciones del diagrama HR, lo que impide diferenciarlas con técnicas observacionales convencionales. Este trabajo evalúa si la astrosismología puede discriminar entre ambas historias evolutivas, en el contexto de la propuesta de misión espacial HAYDN de la ESA.

Se ha simulado la formación de una Blue Straggler mediante transferencia de masa empleando el código MESA, a partir de un sistema binario con una donante de $2.0\,M_\odot$ y una acretora de $1.0\,M_\odot$. El resultado es una estrella rejuvenecida de $1.82136\,M_\odot$ con un perfil de composición química distinto al de una estrella simple. Las frecuencias de oscilación (modos $l=0,1,2$; $n_{pg} \in [-10,10]$) de esta Blue Straggler y de estrellas de referencia de masa similar se han calculado con el código GYRE, comparándolas en 7 puntos significativos de la secuencia principal.

Los resultados muestran diferencias entre ambos escenarios. Los indicadores más sensibles son los modos no radiales ($l\neq0$): en la gran separación ($\Delta\nu$), los modos $l=2$ alcanzan diferencias de hasta ${\sim}\,2.5\,\text{ciclos\,día}^{-1}$ al final de la secuencia principal. La pequeña separación ($\delta\nu_{02}$) y las frecuencias individuales de los modos g muestran diferencias del mismo orden. Se estima un tiempo de integración mínimo de ${\sim}\,142$ días para resolver las diferencias más pequeñas, compatible con las campañas de observación previstas por HAYDN.

El principal obstáculo para la aplicabilidad observacional es la identificación inequívoca de los modos de oscilación en estrellas $\delta$~Scuti, región del diagrama HR donde se ubican las Blue Stragglers, dada la complejidad de su espectro de frecuencias.

\bigskip
\noindent\textbf{Palabras clave:} astrosismología, Blue Stragglers, oscilaciones estelares, MESA, GYRE, transferencia de masa binaria, gran separación, pequeña separación, HAYDN.
